% Kamp 1968, dense-order observation and Chapter IV, Definition 2/Theorem 2.

\subsection{Dense-order separation and a unary basis on the integers}
\label{sec:p09-kamp-iv2-integer-basis}

\subsubsection{The adjacent dense-order observation}

Kamp extends his real-order separation to every particular dense linear
order with more than one point. Such an order may have endpoints, need
not be complete, and need not resemble the reals in cardinality. The
following local argument supplies the extension at those exact hypotheses
\cite[ch.~IV, pp.~99--100]{Kamp1968}.

\begin{hbdefinition}{def:p09-kamp-dense-local-language}{Unary tenses on a fixed dense order}
Fix a strict linear order \(\mathcal J=\langle T,<\rangle\).
A permitted unary \(\mathcal J\)-tense is an operation
\(O:\mathcal P(T)\to\mathcal P(T)\) defined by one source
first-order formula in the equality-free signature \(\{<,P_0\}\)
whose only free variable is the current point
\(t_0\), which contains \(P_0\), and which contains no other
unary predicate or named point. Negation and conjunction act on sets by
complement in \(T\) and intersection. Write
\(\mathsf{TM}_{\mathcal J}\) for the finite formulas generated
from proposition letters \(q_i\) by these operations, with every
permitted unary tense available. As before, every finite
truth-functional connective can be admitted as an abbreviation.

For a convex subset \(I\subseteq T\), complements of subsets of
\(I\) in the next lemma are taken relative to \(I\).
A two-sided sequence of markers in \(I\) is a family
\((d_n)_{n\in\mathbb Z}\) with \(d_n<d_{n+1}\) such that
every \(x\in I\) lies strictly between some two markers.
In particular the markers are cofinal in both directions inside \(I\).
The interpretation of a unary predicate outside \(I\) is unrestricted.
\end{hbdefinition}

\begin{hblemma}{thm:p09-kamp-dense-local-unary}{proved=annex:p9.kamp.dense.local.unary}{Local four-choice closure without completeness}
Let \(T\) be densely ordered, and let \(I\) be a nonempty
convex subset with no first or last point. Let \(A\subseteq T\).
Suppose that \(A\cap I\) is \(\varnothing\), \(I\), a
two-sided marker set \(D=\{d_n:n\in\mathbb Z\}\), or
\(I\setminus D\). Then every permitted unary tense satisfies
\[
 O(A)\cap I\in
       \{\varnothing,I,A\cap I,I\setminus(A\cap I)\}.
\]
More precisely, two points of \(I\) with the same membership in
\(A\) satisfy exactly the same first-order formulas in the signature
\(\{<,P_0\}\) with only the current point free. The formula may
quantify over all of \(T\), including its endpoints and the
unrestricted interpretation outside \(I\).
\end{hblemma}

\begin{hbtheorem}{thm:p09-kamp-dense-unary-separation}{proved=annex:p9.kamp.dense.unary.separation}{Unary separation on every nontrivial dense order}
For every particular dense strict linear order
\(\mathcal J=\langle T,<\rangle\) with more than one element,
neither of the strict operations
\[
\begin{aligned}
 S(A,B)&=\{t\in T:\exists s<t\ (s\in A\land(s,t)\subseteq B)\},\\
 U(A,B)&=\{t\in T:\exists s>t\ (s\in A\land(t,s)\subseteq B)\}
\end{aligned}
\]
is expressible by a finite formula in
\(\mathsf{TM}_{\mathcal J}\) using only its two input letters.
Equivalence is required for every \(A,B\subseteq T\).
The unary tenses may be chosen for this particular order; no uniform
definition across a class of orders is imposed as an additional
hypothesis. Density and more than one element are the only order
hypotheses. On a one-point order the two strict operations are constantly
false, so that excluded boundary behaves differently
\cite[ch.~IV, pp.~99--100]{Kamp1968}.
\end{hbtheorem}

\subsubsection{The integer basis and its equality-free definitions}

\begin{hbdefinition}{def:p09-kamp-iv2-basis}{The four unary integer tenses}
For the remainder of this section the time order is
\(\mathsf{In}=\langle\mathbb Z,<\rangle\).
To display the source's equality-free definitions, abbreviate
\(P_0(z)\land\neg P_0(z)\) by \(\bot_{P_0}(z)\).
This is a contradictory formula, not an extra constant in the language.
The defining formulas at \(t_0\) are
\[
\begin{aligned}
 \pi_-(t_0)={}&\exists t_1\bigl(t_1<t_0\land P_0(t_1)\land
       \forall t_2(t_1<t_2<t_0\to\bot_{P_0}(t_2))\bigr),\\
 \pi_+(t_0)={}&\exists t_1\bigl(t_0<t_1\land P_0(t_1)\land
       \forall t_2(t_0<t_2<t_1\to\bot_{P_0}(t_2))\bigr).
\end{aligned}
\]
The isolation formulas are
\[
\begin{aligned}
 \iota_-(t_0)={}&\exists t_1\bigl(t_1<t_0\land P_0(t_1)\land\\
 &\quad\exists t_2\bigl(t_2<t_1\land\neg P_0(t_2)\land
       \forall t_3(t_2<t_3<t_1\to\bot_{P_0}(t_3))\bigr)\land\\
 &\quad\forall t_4(t_1<t_4<t_0\to\neg P_0(t_4))\bigr),\\
 \iota_+(t_0)={}&\exists t_1\bigl(t_0<t_1\land P_0(t_1)\land\\
 &\quad\exists t_2\bigl(t_1<t_2\land\neg P_0(t_2)\land
       \forall t_3(t_1<t_3<t_2\to\bot_{P_0}(t_3))\bigr)\land\\
 &\quad\forall t_4(t_0<t_4<t_1\to\neg P_0(t_4))\bigr).
\end{aligned}
\]
They define respectively \(\mathsf{PAST}_1\),
\(\mathsf{FUTURE}_1\), \(\mathsf{ISOLP}\), and
\(\mathsf{ISOLF}\). The basis is exactly
\[
 \mathcal S=\{\mathsf{ISOLP},\mathsf{ISOLF},
              \mathsf{PAST}_1,\mathsf{FUTURE}_1,
              \mathsf{NOT},\mathsf{AND}\}.
\]
Its sentential language \(\mathsf{TS}\) uses \(P_1,F_1\)
for immediate past and future, and \(S_1,U_1\) for past and
future isolation. Its finite formulas start with proposition letters
\(q_i\) and apply these four unary constructors and the
sentential constructors \(\neg,\land\).
Disjunction is a Boolean abbreviation. Juxtaposition of unary symbols
denotes nesting; for example \(F_1F_1\alpha=F_1(F_1\alpha)\).
The subscripted \(S_1,U_1\) are unary isolation operators;
the unsubscripted \(S,U\) remain the binary target operations
\cite[ch.~IV, def.~2 and thm.~2, p.~100]{Kamp1968}.
\end{hbdefinition}

\begin{hblemma}{thm:p09-kamp-iv2-unary-semantics}{proved=annex:p9.kamp.iv2.unary.semantics}{The exact integer clauses for the source basis}
For any \(\mathsf{TS}\)-formula \(\alpha\), any valuation,
and any \(i\in\mathbb Z\), write \(\alpha(i)\) for its
truth at \(i\). Then
\[
\begin{aligned}
 (P_1\alpha)(i)&\Longleftrightarrow\alpha(i-1),\\
 (F_1\alpha)(i)&\Longleftrightarrow\alpha(i+1),\\
 (S_1\alpha)(i)&\Longleftrightarrow
   \exists j<i\bigl(\alpha(j)\land\neg\alpha(j-1)\land
       \forall k(j<k<i\to\neg\alpha(k))\bigr),\\
 (U_1\alpha)(i)&\Longleftrightarrow
   \exists j>i\bigl(\alpha(j)\land\neg\alpha(j+1)\land
       \forall k(i<k<j\to\neg\alpha(k))\bigr).
\end{aligned}
\]
Thus an isolation operator requires the nearest strict occurrence
in its direction to have a false neighbor on its far side.
There is no condition on \(\alpha(i)\). The intervening
universal condition is vacuous when that nearest occurrence is an
immediate neighbor
\cite[ch.~IV, pp.~101--102]{Kamp1968}.
\end{hblemma}

\subsubsection{The complete two-input construction}

\begin{hbdefinition}{def:p09-kamp-iv2-marker-formulas}{Markers that distinguish a successful endpoint from a failure}
Use \(q_0\) for the witness proposition and \(q_1\) for
the open-interval proposition. Define the following formulas of
\(\mathsf{TS}\):
\[
\begin{aligned}
 C&=\neg q_0\land q_1,\\
 B_P&=\neg q_0\land\neg q_1\land P_1\neg q_0
                                      \land P_1P_1\neg q_0,\\
 B_F&=\neg q_0\land\neg q_1\land F_1\neg q_0
                                      \land F_1F_1\neg q_0,\\
 \eta_1&=(q_0\land P_1\neg q_0\land P_1q_1)
                                      \lor B_P\lor F_1B_P,\\
 \eta_2&=(q_0\land F_1\neg q_0\land F_1q_1)
                                      \lor B_F\lor P_1B_F.
\end{aligned}
\]
The marker \(\eta_1\) is used for a future search and inspects
the two preceding positions. The marker \(\eta_2\) is used
for a past search and inspects the two following positions.
The final shifted disjunct places an extra marker immediately before
a failed future search, or immediately after a failed past search.

For \(\delta\in\{+1,-1\}\), define the directional notation
\[
\begin{array}{c|ccc}
 \delta&N_\delta&J_\delta&\eta_\delta\\ \hline
 +1&F_1&U_1&\eta_1\\
 -1&P_1&S_1&\eta_2
\end{array}
\]
where \(\eta_\delta\) is only a directional abbreviation in
this table: \(\eta_{+1}=\eta_1\) and
\(\eta_{-1}=\eta_2\). The symbol \(N_\delta\) means one
step in the chosen direction; \(N_{-\delta}\) means one step
back toward the starting point.
\end{hbdefinition}

\begin{hblemma}{thm:p09-kamp-iv2-marker-lemma}{proved=annex:p9.kamp.iv2.marker.lemma}{The marker calculation in both directions}
Fix any integer \(t\), either \(\delta\in\{+1,-1\}\),
and any interpretations of \(q_0,q_1\). Suppose that \(C\)
holds at \(t+\delta\) and \(t+2\delta\). Then
\[
 (J_\delta\eta_\delta)(t)
 \quad\Longleftrightarrow\quad
 \exists m\geq3\bigl(q_0(t+\delta m)\land
       \forall r\,(1\leq r<m\to q_1(t+\delta r))\bigr).
\]
Here \(m,r\) range over integers.
The proof includes a successful first endpoint, a first failed
interval point, and the case with no endpoint or failure anywhere
in the chosen direction. It also verifies the source's alternative
case in which the witness proposition never occurs in that direction.
\end{hblemma}

\begin{hbdefinition}{def:p09-kamp-iv2-target-formulas}{Explicit expressions for strict UNTIL and SINCE}
Retain both conjuncts of the source construction. Put
\[
\begin{aligned}
 G_F&=F_1q_0\lor
       \bigl(F_1q_1\land(F_1F_1q_0\lor F_1F_1q_1)\bigr),\\
 G_P&=P_1q_0\lor
       \bigl(P_1q_1\land(P_1P_1q_0\lor P_1P_1q_1)\bigr),\\
 \zeta_1&=G_F\land\bigl[
       F_1q_0\lor(F_1q_1\land F_1F_1q_0)\\
       &\hspace{38mm}\lor
          (F_1C\land F_1F_1C\land U_1\eta_1)\bigr],\\
 \zeta_2&=G_P\land\bigl[
       P_1q_0\lor(P_1q_1\land P_1P_1q_0)\\
       &\hspace{38mm}\lor
          (P_1C\land P_1P_1C\land S_1\eta_2)\bigr].
\end{aligned}
\]
All abbreviations expand to finite formulas over exactly
\(\mathcal S\), with no proposition letters beyond \(q_0,q_1\).
These give the corrected source construction: the marker inspections
face toward the starting point, and the repeated shifts in each
outer formula face in its search direction
\cite[ch.~IV, thm.~2, pp.~100--105]{Kamp1968}.
\end{hbdefinition}

\begin{hbtheorem}{thm:p09-kamp-iv2}{proved=annex:p9.kamp.iv2}{Kamp's integer-order unary basis theorem}
For every \(A,B\subseteq\mathbb Z\), every
\(t\in\mathbb Z\), and the valuation \(q_0\mapsto A\),
\(q_1\mapsto B\),
\[
\begin{aligned}
 \zeta_1(t)&\Longleftrightarrow
       \exists s>t\bigl(s\in A\land
                       \forall u(t<u<s\to u\in B)\bigr),\\
 \zeta_2(t)&\Longleftrightarrow
       \exists s<t\bigl(s\in A\land
                       \forall u(s<u<t\to u\in B)\bigr).
\end{aligned}
\]
Consequently both strict UNTIL and strict SINCE are expressible
on \(\mathsf{In}\) from
\(\{\mathsf{ISOLP},\mathsf{ISOLF},\mathsf{PAST}_1,
\mathsf{FUTURE}_1,\mathsf{NOT},\mathsf{AND}\}\).
The first operand remains the witness proposition and the second
the open-interval proposition. An immediate witness makes the
interval requirement vacuous. The conclusion holds at every integer,
including when either predicate is empty, full, or unbounded in
one or both directions
\cite[ch.~IV, thm.~2, pp.~100--105]{Kamp1968}.
\end{hbtheorem}

\begin{hbexample}{def:p09-kamp-iv2-examples}{Successful endpoints, paired failures and immediate witnesses}
At \(t=0\), take \(A=\{3,4,5,\ldots\}\) and
\(B=\mathbb Z\). The first two future points satisfy
\(C\). The marker \(\eta_1\) is true at three and
false at every other positive integer. Thus \(U_1\eta_1\)
and \(\zeta_1\) are true, with three a witness for strict
UNTIL. Reflecting these sets through zero gives the corresponding
SINCE example with \(A=\{\ldots,-5,-4,-3\}\).

If instead \(A=\varnothing\) and \(B=\{1,2\}\), the
first failed future point is three. The marker is true at both two
and three, so its nearest future occurrence is not isolated on its
far side. Thus \(U_1\eta_1\) and \(\zeta_1\) are false.
If \(A=\varnothing\) and \(B=\mathbb Z\), there is no
future marker at all, and the same conclusion follows.

For an immediate witness take \(A=\{1\}\) and any \(B\),
even \(B=\varnothing\). Then \(F_1q_0\) makes both
conjuncts of \(\zeta_1\) true, matching the empty interval
\((0,1)\cap\mathbb Z\). For \(A=\{2\}\),
\(\zeta_1(0)\) is true exactly when \(1\in B\).
Neither case imposes a condition on membership of the current point
or the witness in \(B\). The annex verifies the examples and
their reflected past cases.
For comparison, \(A=\{0\}\) and \(B=\{-1,-2\}\)
make \(\zeta_2(0)\) false: the current point cannot be
used as a strict past witness.
\end{hbexample}
